% Generated draft: requires mathematical and source review.
\documentclass[11pt]{article}
\usepackage[margin=25mm]{geometry}
\usepackage{fontspec}
\setmainfont{Latin Modern Roman}
\usepackage{amsmath,amssymb,mathtools,unicode-math}
\setmathfont{Latin Modern Math}
\usepackage{enumitem,needspace,xurl,hyperref}
\setlength{\parindent}{0pt}
\setlength{\parskip}{5pt}
\setlength{\emergencystretch}{3em}
\newcommand{\N}{\mathbb{N}}
\newcommand{\Z}{\mathbb{Z}}
\newcommand{\Q}{\mathbb{Q}}
\newcommand{\R}{\mathbb{R}}
\newcommand{\C}{\mathbb{C}}
\newcommand{\F}{\mathbb{F}}
\newcommand{\A}{\mathbb{A}}
\newcommand{\PP}{\mathbb{P}}
\newcommand{\E}{\mathbb{E}}
\newcommand{\eps}{\varepsilon}
\newcommand{\dd}{\,\mathrm{d}}
\begin{document}
\section*{UnsolvedMath ID: 30003213}
\textbf{Source problem code:} catalogue entry
\subsection{Definitions and mathematical statement}
The supplied record presents a question about the strong Atiyah conjecture, but its dated literature assessment says that the underlying extract is an opening fragment without a complete question. These conflicting descriptions are unresolved. The formulation below is therefore a conditional reconstruction of the displayed registry question.

Let \(F\) be a finitely generated free group, let \(N\triangleleft F\), and write \(G=F/N\), with quotient map \(q:F\to G\). Bounded finite-subgroup orders means that there is an integer \(B\geq 1\) such that \(|H|\leq B\) for every finite subgroup \(H\leq G\). Under this hypothesis define
\[
L=\operatorname{lcm}(G):=\operatorname{lcm}\{\,|H|:H\leq G\text{ is finite}\,\}.
\]
This is a finite positive integer; the trivial subgroup is included.

The source specifies only \(K\subset\C\). To interpret \(K[F]\) as a group algebra, provisionally assume that \(K\) is a unital subring of \(\C\). Whether a field or some other coefficient class was intended is not determined by the supplied evidence. With this provisional assumption, \(K[F]\) consists of finite sums \(\sum_{f\in F}c_f f\), with coefficients in \(K\), and \(q\) induces a homomorphism \(q_*:K[F]\to K[G]\).

Here is one explicit operator convention, since the source does not define \(\phi_G^A\). For positive integers \(n,m\) and \(A\in\operatorname{Mat}_{n\times m}(K[F])\), write
\[
(q_*A)_{ij}=\sum_{g\in G}a_{ij,g}g.
\]
Let \(\ell^2(G)\) be the Hilbert space of square-summable complex functions on \(G\), with standard basis \((\delta_h)_{h\in G}\). Define the bounded operator
\[
\phi_G^A:\ell^2(G)^n\longrightarrow\ell^2(G)^m,
\qquad
\phi_G^A(\delta_h e_i)=\sum_{j=1}^{m}\sum_{g\in G}a_{ij,g}\delta_{hg}e_j,
\]
where \(e_i\) and \(e_j\) denote coordinate vectors in the respective spaces. This operator commutes with the diagonal left translation action of \(G\). If \(P_A\) is the orthogonal projection onto its kernel, define its von Neumann dimension by
\[
\dim_G\ker\phi_G^A
:=\sum_{i=1}^{n}\left\langle P_A(\delta_{1_G}e_i),\delta_{1_G}e_i\right\rangle.
\]

Under these explicitly provisional coefficient and operator conventions, the proposed universal question is: for every finitely generated free group \(F\), every normal subgroup \(N\triangleleft F\) such that \(G=F/N\) has bounded finite-subgroup orders, every unital subring \(K\subset\C\), all positive integers \(n,m\), and every \(A\in\operatorname{Mat}_{n\times m}(K[F])\), must
\[
\dim_G\ker\phi_G^A\in\frac{1}{\operatorname{lcm}(G)}\Z
\]
hold? The coefficient class, matrix-size convention, operator convention, and correspondence of this completed question with the report require human verification.
\subsection{Short English statement}
The displayed registry question asks whether kernel dimensions of matrices over a free-group algebra, acting through a quotient group, have denominators dividing the least common multiple of the orders of its finite subgroups. The formulation is conditional because the coefficient domain and operator are unspecified, and the supplied dated assessment disputes whether the source extract contains a complete question. The supplied \emph{open} label is not independently verified.
\subsection{Sources}
\begin{itemize}
\item Supplied local registry record: \texttt{data/unsolvedmath/problems.json}, permanent UnsolvedMath ID \(30003213\), title \emph{Strong Atiyah Conjecture for Free-Group Quotients}, source problem code \texttt{catalogue entry}, in the \emph{Oberwolfach Reports Open Problems} collection. The proposed year is recorded as \(2016\); no individual proposer is supplied.
\item Supplied report citation: \emph{Measured Group Theory (2017)}, Oberwolfach Reports, pp.\ 2347--2397, DOI \texttt{10.4171/owr/2016/41}. Supplied source URL and literature-source link, labeled ``source report'': \url{https://doi.org/10.4171/owr/2016/41}. The background describes Oberwolfach Reports as documenting mathematical workshops and their open-problem discussions. This report has not been independently consulted for this draft.
\item Literature assessment dated \texttt{2026-08-21}: supplied dated curation, not independently verified. It labels the record open, describes the extracted text as only an opening fragment of a discussion of the strong Atiyah conjecture, and says that no definite proposition or complete question, and hence no stronger status, can be justified from that exact extract. Its listed source checked is the report linked above. The background expressly calls for independent verification of time-sensitive status and cited claims.
\item The registry separately records \texttt{statement\_status: corrected\_verified} and asserts that the clean statement was checked against the official report and preserves its hypotheses, notation, and open target. These are supplied verification claims; their conflict with the dated fragment assessment remains unresolved.
\end{itemize}
\end{document}
